\subsection{5.4 网页与微信小程序共用服务的设计}
\noindent 两个客户端只承担各自的表示层；服务端按内部顾客 ID 调用同一库存、购物车、预约、合同和结算领域规则。门店员工继续使用网页工作台。微信扫码入口、微信登录与 BSR 账号绑定属于渠道适配，不得把微信 OpenID、网页 Cookie 或二维码场景当业务订单主键。

\subsubsection{5.4.1 分层责任与接口边界}
\begin{BSRTable}
\rowcolor{TableHead}\color{white}\bfseries 层 / 模块 & \color{white}\bfseries 责任、数据及不可越界的规则 \\
网页与小程序表示层 & 网页顾客和微信小程序均提供日期段、分类/尺码/款式筛选、购物车、规则阅读、报价及本人订单；员工只在网页端执行备货、交付、检修和复核。两端按状态文字、责任人和更新时间显示进度，不将本地缓存当成交付事实。\\
身份与渠道适配 & 网页使用受保护的 Cookie 与 Origin 校验；小程序用一次性微信临时码、现有账号证明和独立 Bearer 会话。适配器只将已验证的微信身份映射到顾客 ID；停用、改密、解绑、退出及恢复后旧令牌失效。\\
领域服务 & 共享 Cart、Quote、Reservation、InventoryUnit、Contract、DamageCase 和 DepositSettlement 规则；小程序不能绕过店员/经理的对象授权，也不能因另一端缓存较新就跳过服务器重核。\\
数据访问 & wechat\_bindings、mini\_login\_challenges、mini\_code\_uses、mini\_sessions、mini\_quotes 保存身份和短时报价；customer\_carts 及行、订单和单件库存仍由两个客户端共用。身份表可迁移，不建立小程序专用订单或库存副本。\\
外部依赖 & 正式 AppID 已核对；微信身份交换和正式小程序码由服务端调用平台，凭据留在私有配置；HTTPS 合法请求域名、平台审核与手机扫码仍待验。模拟收退和签署不能对顾客宣称为真实交易。\\
\end{BSRTable}

\subsubsection{5.4.2 事务、数据保护和恢复}
\noindent 服务端先以账号证明明确绑定，活跃 AppID/OpenID 摘要与顾客均设唯一约束；临时码、挑战和会话都有失效时刻及一次性/撤销标记。库中只存 OpenID 的密钥化摘要与令牌摘要，原始 OpenID、微信会话密钥、AppSecret、临时码、Bearer 不进入日志、公开静态目录或 Git 发布包。解绑、账号停用和恢复副本切换须撤销旧小程序会话；网页会话按自身机制分别处理。

\noindent 共享购物车用 version 检测另一端更改；报价保存客户、购物车版本、NOK 金额指纹和短时到期。预约在一事务内核对报价指纹、条款版本、顾客水平、日期段单件可用量和请求幂等引用。冲突只返回重读或重报，不提交半张订单。两端订单列表、详情和具名进度都以服务端同一预约 ID 为准。SSE 仅帮助网页重新查询，不能作为小程序授权或库存事实。

\subsubsection{5.4.3 当前证据与正式验收边界}
\begin{BSRTable}
\rowcolor{TableHead}\color{white}\bfseries 能力 & \color{white}\bfseries 本机状态 / 仍需验证 \\
共享业务 & 隔离合成 SHOWCASE 已验证共享购物车、最后一件并发、过期报价、同键重试及同一订单经网页员工备货、交付、正常归还/受损鉴定和模拟结算；本人原图受控读取与退出清理已有测试，真机预览待验。\\
原生顾客端 & 正式 AppID 已在开发者工具核对，模拟器目录、购物车、订单与备案信息可见，所查页面问题数为零。本地 local-demo 只连回环演示；真实微信身份、扫码与发布仍待验。\\
网站入口 & 未取得正式平台码前保持关闭/筹备；不得展示示意二维码并暗示可扫码。已取得码后须核对 AppID、签发者、场景、有效期和停用回退，再进行 10 次手机扫码验收。\\
资金与签署 & 本机状态均标模拟；正式支付、原路退款、电子签署、对账及监管要求另行集成和验收。8\% 押金及超时/受损费率仍为待 BSR 批准的演示政策。\\
\end{BSRTable}
